\documentclass[11pt]{article}

\usepackage{amsmath,amssymb,amsthm,mathtools}
\usepackage[margin=1in]{geometry}
\usepackage{url}

\newtheorem{theorem}{Theorem}[section]
\newtheorem{lemma}[theorem]{Lemma}
\newtheorem{corollary}[theorem]{Corollary}
\newtheorem{remark}[theorem]{Remark}

\newcommand{\Gap}{\Gamma_H}

\begin{document}

\section{Purpose}

This note strengthens the high-density comparison between the constant
graphon and balanced complete multipartite graphons.  Instead of testing one
fixed endpoint, it treats every bounded normalized complement direction
\[
 W_q=1-qU,
 \qquad \int U=1,
 \qquad 0\leq U\leq M.
\]
For every graph containing a triangle, the chromatic lower-bound gap is
strictly positive for all sufficiently small $q>0$, uniformly over every
such $U$ with the fixed bound $M$.  An explicit sufficient range is given.

The theorem is asymptotic in $p=1-q\to1$ and does not assert global
positivity.  It is nevertheless a full infinite-dimensional test: $U$ is an
arbitrary bounded measurable symmetric kernel, not a step graphon or a
selected rewiring direction.

\section{Cycle-rank expansion}

Let $H$ be a finite simple graph with $n$ vertices and $m$ edges.  For
$A\subseteq E(H)$, let $H_A=(V(H),A)$, let
\[
 r(A)=n-c(H_A),\qquad \beta(A)=|A|-r(A),
\]
where isolated vertices count as components.  Thus $r$ is the graphic-matroid
rank and $\beta$ is its nullity.  Define
\[
 \Gamma_H(W)=t(H,W)-\Phi_H(p(W)),
 \qquad
 \Phi_H(p)=(1-p)^n\chi_H\!\left(\frac1{1-p}\right).
\]

\begin{lemma}[Exact normalized-complement identity]
Let $U$ be a bounded nonnegative symmetric kernel on a probability space
with $\int U=1$.  If $W_q=1-qU$ is a graphon, then
\begin{equation}
 \Gamma_H(W_q)
 =\sum_{A\subseteq E(H)}(-1)^{|A|}q^{r(A)}
 \left(q^{\beta(A)}t(H_A,U)-1\right).
 \label{eq:cycle-rank}
\end{equation}
\end{lemma}

\begin{proof}
Expanding one factor $1-qU$ on every edge gives
\[
 t(H,1-qU)=
 \sum_{A\subseteq E(H)}(-1)^{|A|}q^{|A|}t(H_A,U).
\]
Whitney's rank expansion gives
\[
 q^n\chi_H(1/q)
 =\sum_{A\subseteq E(H)}(-1)^{|A|}q^{r(A)}.
\]
Since $p(W_q)=1-q$, subtracting the second identity from the first and using
$|A|=r(A)+\beta(A)$ proves \eqref{eq:cycle-rank}.  Both sums are finite, so
no interchange or convergence argument is involved.
\end{proof}

Let
\[
 P_2(H)=\sum_{v\in V(H)}\binom{d_H(v)}2
\]
be the number of unordered adjacent edge pairs, and let $N_3(H)$ be the
number of unoriented triangles.  If
$d_U(x)=\int U(x,y)\,dy$, then
\[
 t(P_3,U)=\int d_U(x)^2\,dx\geq
 \left(\int d_U\right)^2=1.
 \label{eq:path-cauchy}
\]

\begin{lemma}[Complete quadratic coefficient]
For every fixed bounded $U$ as above,
\begin{equation}
 \Gamma_H(1-qU)
 =q^2\left(
 P_2(H)(t(P_3,U)-1)+N_3(H)
 \right)+O_{H,U}(q^3).
 \label{eq:quadratic}
\end{equation}
\end{lemma}

\begin{proof}
All rank-zero and rank-one terms in \eqref{eq:cycle-rank} vanish because
$t(\varnothing,U)=1$ and $t(K_2,U)=\int U=1$.  A rank-two forest with two
edges is either $P_3$ or $2K_2$.  The latter has density one, while each
adjacent edge pair contributes $t(P_3,U)-1$.  The only cyclic simple graph
of rank two is a triangle.  Its contribution is
\[
 -q^2(qt(K_3,U)-1)=q^2-q^3t(K_3,U).
\]
Every other edge set has rank at least three.  Collecting the displayed
terms proves \eqref{eq:quadratic}.
\end{proof}

The strict triangle contribution is important.  Degree irregularity of $U$
can only add the nonnegative term in \eqref{eq:path-cauchy}; regularity is
not assumed.

\section{An explicit uniform remainder}

For $M\geq1$, define
\begin{equation}
 C_H(M)=
 \sum_{j=3}^{m}\binom mj M^j
 +\sum_{j=3}^{m}\binom mj-N_3(H).
 \label{eq:remainder-constant}
\end{equation}
The first sum bounds all cyclic and rank-at-least-three density terms.  The
second counts the rank-at-least-three edge sets: every set with at least four
edges has rank at least three in a simple graph, while among three-edge sets
exactly the triangles have rank two.

\begin{theorem}[Bounded-complement high-density stability]
Let $H$ contain a triangle, let $M\geq1$, and let $U$ be any measurable
symmetric kernel satisfying
\[
 0\leq U\leq M,\qquad \int U=1.
\]
If
\begin{equation}
 0<q<\min\left\{\frac1M,
 \frac{N_3(H)}{C_H(M)}\right\},
 \label{eq:q-range}
\end{equation}
then $W_q=1-qU$ is a graphon of edge density $1-q$ and
\[
 \Gamma_H(W_q)>0.
\]
More precisely,
\begin{equation}
 \Gamma_H(W_q)
 \geq q^2\bigl(N_3(H)-qC_H(M)\bigr)>0.
 \label{eq:explicit-lower}
\end{equation}
\end{theorem}

\begin{proof}
The condition $qM<1$ ensures $0\leq W_q\leq1$.  The quadratic forest term
is nonnegative by \eqref{eq:path-cauchy}.  The order-three part of each
triangle has absolute value at most $q^3M^3$.  For every $A$ of rank at
least three, $q\leq1$ and $t(H_A,U)\leq M^{|A|}$ give
\[
 \left|q^{r(A)}
 \left(q^{\beta(A)}t(H_A,U)-1\right)\right|
 \leq q^3\left(M^{|A|}+1\right).
\]
There are $\binom mj$ edge subsets of size $j$.  Separating the
$N_3(H)$ triangles from the rank-at-least-three sets yields exactly
\eqref{eq:remainder-constant}.  Subtracting this absolute remainder from the
quadratic triangle contribution proves \eqref{eq:explicit-lower}, and
\eqref{eq:q-range} makes it strict.
\end{proof}

\begin{corollary}[All current open rows]
Every one of the current $29$ open Atlas rows satisfies the conclusion of
the theorem for each fixed $M\geq1$.  Their Atlas IDs are
\[
\begin{gathered}
43,118,122,124,126,127,130,137,139,145,147,148,151,152,153,\\
157,160,168,169,171,174,178,181,185,188,194,196,199,203.
\end{gathered}
\]
Their triangle counts range from one to nine.
\end{corollary}

\begin{proof}
Exact finite enumeration verifies that every displayed graph contains a
triangle.  The verifier independently counts adjacent edge pairs and
triangles and reconstructs $C_H(M)$ for every row.
\end{proof}

\section{Relation to the two endpoint comparisons}

The constant graphon is $U\equiv1$, hence is covered with $M=1$.  The exact
constant-versus-target polynomial in the companion note is stronger for that
one direction because it proves the correct sign on the entire admissible
density interval.

The normalized complement of the balanced complete $k$-partite graphon has
value $k$ on its diagonal blocks and zero elsewhere.  Along
$q=1/k\to0$, its bound is $M=k=1/q$, not a fixed constant.  It therefore
lies exactly outside the uniform regime of this theorem, as it must: those
graphons attain equality.  Any unresolved high-density construction must
similarly concentrate its complement so that
$\|(1-W)/(1-p)\|_\infty$ diverges.  A fixed bounded complement shape cannot
be negative for a triangle-containing graph.

\section{Limitations and reproducibility}

The theorem requires $N_3(H)>0$ and a fixed $L^\infty$ bound $M$.  It does
not cover even-girth graphs, whose first cyclic term can have the opposite
sign, or $q$-dependent normalized complements with unbounded amplitude.  It
does not classify any row by itself.

Run
\begin{verbatim}
python codes/principled_negative_tests.py
\end{verbatim}
from the repository root.  The bounded-complement audit uses only exact
finite graph enumeration and integer polynomials; it reports all $29$ rows
and the triangle-count range $1\ldots9$.

\end{document}
